%%
% ICICS / Springer LNCS（本目录 llncs.cls + splncs04.bst）。
% 正文自 paperzh short.tex 迁入并按 samplepaper.tex 版式调整（runningheads、无 geometry）。
% 编译（在 ICICS模版 目录下）：
%   pdflatex ciphernova_icics
%   bibtex ciphernova_icics
%   pdflatex ciphernova_icics
%   pdflatex ciphernova_icics
% 依赖：本目录 refs.bib（与仓库根目录 refs.bib 请自行保持同步）；插图等见 graphicspath。
\documentclass[runningheads]{llncs}

\usepackage[UTF8]{ctex}
\usepackage{graphicx}
\graphicspath{{../figures/}{../}{figures/}}
\usepackage{float}
\usepackage{amsmath,amssymb}
\usepackage[super,square,numbers,sort&compress]{natbib}
\usepackage{enumitem}
\setlist[enumerate]{leftmargin=1.2em, itemindent=0em, labelwidth=1em, labelsep=0.4em}
\setlist[itemize]{leftmargin=1.2em, labelsep=0.4em}
\usepackage[table]{xcolor}
\usepackage{booktabs}
\usepackage{array}
\usepackage{multirow}
\usepackage{hyperref}
\hypersetup{colorlinks=true, linkcolor=blue, citecolor=blue, urlcolor=blue}
\usepackage{caption}
\captionsetup[figure]{aboveskip=4pt, belowskip=4pt}
\captionsetup[table]{aboveskip=4pt, belowskip=4pt}
\setlength{\textfloatsep}{6pt plus 2pt minus 1pt}
\setlength{\intextsep}{6pt plus 2pt minus 1pt}

\begin{document}
\title{CIPHERNOVA：一种基于测试反馈与分层提示的密码学代码生成方法}
\titlerunning{CIPHERNOVA：测试反馈与分层提示驱动的密码学代码生成}

% ===== Anonymous switch =====
\newif\ifanonymized
\anonymizedtrue      % 投稿匿名版
% \anonymizedfalse   % 终稿实名版

\ifanonymized
\author{Anonymous Authors}
\institute{Anonymous Institution\\
\email{anonymous@anonymous.edu}}
\authorrunning{Anonymous Authors}
\else
\author{Jia-Yi He\inst{1} \and Xiao Lan\thanks{Corresponding author: lanxiao@scu.edu.cn}}
\institute{School of Cyber Science and Engineering, Sichuan University, Chengdu, Sichuan, China\\
\email{2022141500268@stu.scu.edu.cn, lanxiao@scu.edu.cn}}
\authorrunning{J.\ He and X.\ Lan}
\fi
\maketitle

\begin{abstract}
传统大语言模型（LLM）在密码学代码生成中普遍面临错误率高、缺乏领域优化等问题。本文提出系统化方法：（1）基于测试反馈的自适应改进——闭环验证、标准向量、错误模式识别与针对性修复；（2）分层可组合提示——四层与 101 模板及；\textbf{反复失败驱动的提示优化}；（3）密码学错误分类与修复策略；（4）兼容离线本地模型（如 Qwen2.5-Coder:7B）；（5）\textbf{教师示例检索注入}——以云端通过实现为教师，在生成/改进时注入少样本与同题参考，不重训练即可提升本地模型通过率。在 3095 例标准向量网格评测下，云端 LLM 达 100\% 生成提取成功率与 100\% 功能测试通过率；本地qwen2.5coder亦达88.89\% FTPR。方法在同口径上显著优于无提示基线，结论限于本文基准与评测设置。
\end{abstract}

\keywords{密码学代码生成 \and 大语言模型 \and 测试反馈 \and 提示工程 \and 错误模式分类 \and 代码自改进}

\section{引言}
\subsection{研究背景与问题陈述}
密码算法是保密通信、身份认证与数据隐私的常见基座，其实现既要覆盖分组密码、公钥密码等不同形态，又须满足 NIST、FIPS 等规范条款；填充、初始向量等细节一旦处理不当，便可能破坏机密性或完整性。

非专长开发者依赖文档与开源库时常面临学习成本高、定制难等风险。GitHub Copilot、CodeT5~\cite{codet5} 及系列代码大模型~\cite{codex,codellama,deepseekcoder,starcoder} 推动了编程辅助，但在密码学任务上仍突出表现为：细节敏感导致功能失败率高；通用提示难以覆盖标准符合性与安全约束；失败依赖人工纠错；Pass@k 等通用指标~\cite{codex} 难充分反映密码代码的功能与安全要求。

因此，如何面向密码学领域为 LLM 设计代码生成优化方法，以提升代码的标准符合性，已成为亟待解决的重要问题。

\subsection{错误密码代码的安全影响}
\label{subsec:security_implications}
错误或不符合标准的实现会削弱部署安全性；CVE 与产业案例表明，弱 IV/ECB 误用、弱 KDF 与填充错误等可导致机密性或完整性失效~\cite{secsurvey,cryptollm}。降低此类代码的生成与部署比例具有\emph{安全保证}意义：本文将标准从\textbf{可编译}提升为\textbf{通过标准向量并满足密码学约束}，以减少 LLM 输出在安全敏感场景中的隐患。

\subsection{研究目标与主要贡献}
针对上述问题，本文提出一种基于测试反馈与分层提示的密码学代码生成方法。核心目标是通过测试反馈的闭环改进与分层提示的精准引导，提升 LLM 生成的密码学代码质量，减轻开发者的负担。

本文的主要贡献可概括为：
\begin{itemize}
    \item 提出\textbf{基于密码学代码错误模式分类的测试反馈与自适应改进机制}。建立四类错误的分类框架并实现六步闭环流程；自动识别错误类型并生成针对性修复指导。
    \item 设计\textbf{组合领域提示工程方法}。构建四层提示架构，通过迭代强化提示以修复反复出现的生成失败，形成精炼模板。
    \item 研制引入\textbf{云端LLM教师示例检索注入}面向\textbf{离线本地大模型}（如 Qwen2.5-Coder:7B）的生成引擎，在无公网环境下仍可部署；以通过标准测试的云端实现为教师，在生成与改进阶段注入少样本及同题参考，\textbf{无需重训练}即可提升本地模型的功能测试通过率，兼顾数据隐私。
\end{itemize}
实验构建了包含 3095 个真实测试用例的密码学代码生成数据集，覆盖 4 种主流算法与 3 种编程语言，为相关研究提供了标准化评估基准。通过对比实验与消融实验，验证了所提方法的有效性及各模块的贡献。


\subsection{论文组织结构}
本文后续章节安排如下：第 2 节综述相关研究；第 3 节阐述方法的整体架构及三个核心创新模块的设计细节；第 4 节通过实验验证方法性能；第 5 节总结全文并讨论局限与未来方向。

\section{相关工作}
本节从代码 LLM、领域代码生成、提示工程与自改进四个方面进行综述。

\subsection{大语言模型代码生成}
LLM 代码生成依赖大规模预训练与自回归建模。CodeGen/ CodeGen2~\cite{codegen1,codegen2}、Codex~\cite{codex}、CodeT5/+~\cite{codet5,codet5plus}、InCoder~\cite{incoder} 及 StarCoder、WizardCoder、Code Llama、DeepSeek-Coder~\cite{starcoder,wizardcoder,codellama,deepseekcoder} 等推动了通用程序合成；CodeBERT/GraphCodeBERT~\cite{codebert,graphcodebert} 与 CodeXGLUE~\cite{codexglue} 强化了表示与评测基础设施。总体面向通用代码，语料中密码学场景稀疏且缺少标准符合性约束。本文以\textbf{云端与本地 LLM}为骨干，\textbf{以领域提示与测试反馈闭环}增强其在标准向量下的密码学代码生成能力。

\subsection{领域代码生成}
领域注入是提升专业代码质量的关键。密码学方面，EncGPT~\cite{encgpt} 以多智能体工作流动态生成加密算法；CryptoLLM~\cite{cryptollm} 探索 LLM 对密码 API 误用的提示式检测；二者尚缺系统错误模式刻画与测试驱动闭环。SecCoder~\cite{seccoder} 表明领域知识可提升安全代码泛化。但在标准符合性、错误分类与“生成—验证—改进”一体化上仍有空白。\textbf{受 EncGPT 的结构化思路与 SecCoder 的注入经验启发，我们将错误类别显式写入提示与修复指导}，形成按算法/语言区分的约束。

\subsection{提示工程}
CoT/零样本 CoT/Plan-and-Solve~\cite{cot,zeroshotcot,planandsolve} 与少样本程序综合~\cite{fewshot,progsynprompt} 提升了推理与生成质量；检索式与仓库级提示~\cite{coderagbench,repolevelprompt} 拓展了上下文组织；分层提示~\cite{hierprompt} 与 Self-Refine~\cite{selfrefine} 强调结构与迭代精炼。现有工作未针对密码学给出可随算法、语言与模型组合适配的分层模板。\textbf{本文将分层可组合提示~\cite{hierprompt} 与迭代精炼~\cite{selfrefine} 实例化为四层架构及失败驱动精炼}，为各组合提供定制约束。

\subsection{代码自改进技术}
“生成—验证—改进”在 Self-Edit~\cite{selfedit}、RefineCoder~\cite{refinecoder}、AgentCoder~\cite{agentcoder}、RepairAgent~\cite{repairagent}、ReCode~\cite{recode} 等工作中得到验证；与之相呼应的还包括基于执行反馈的学习式调试~\cite{cycle,rlef}、结合程序分析的 LLM 修复与分层/多智能体调试~\cite{ledex,li2023,rgd,unidebugger,tang2024}，以及安全代码方向的综述与指令微调~\cite{secsurvey,securecodegen}。总体仍偏通用修复或安全代码，缺少面向密码学向量测试的错误分类与模板化修复。本文\textbf{采用闭环实现六步流程}，由\emph{密码学}错误模式与领域模板驱动，将失败映射为可执行提示。

\subsection{研究空白小结}
\label{subsec:research_gaps}

据我们所知，尚少工作将\emph{测试反馈闭环}、密码学场景的\emph{分层可组合提示}与\emph{系统化错误模式分类}统一于同一框架。空白主要体现在：缺少可操作的密码学错误分类；缺少随算法/语言/LLM 适配的组合式提示；缺少与标准向量紧耦合的自改进流程。本文据此给出一体化方案。

\section{方法设计}
基于上述研究空白，并依托通用代码 LLM 与领域注入~\cite{encgpt,seccoder}、分层/迭代提示~\cite{hierprompt,selfrefine} 及闭环自改进~\cite{selfedit,refinecoder,agentcoder}，本节给出具体方法。整体思路与五层架构见图~\ref{fig:framework}：交互与 API 层之上，\emph{CryptoAgent} 负责提示拼装、调用、提取/清洗与验证，\emph{LLMAdapter} 对接云端与本地 Qwen2.5-Coder:7B；\textbf{模块 A} 在失败时触发改进闭环；工具层实现 \textbf{模块 B}（分层提示）与 \textbf{模块 C}（错误模式分类）；底层为外部模型服务。

\begin{figure}[htbp] 
  \centering
  \includegraphics[width=\textwidth]{Framework.png} 
  \caption{基于测试反馈与分层提示的密码学代码生成系统整体框架}
  \label{fig:framework}
\end{figure}
本方法主要由三个核心模块构成：基于测试反馈的自适应改进机制、分层可组合的领域提示工程方法，以及密码学代码错误模式分类系统。


\subsection{基于测试反馈的自适应改进机制}
基于测试反馈的自适应改进机制是本文方法的核心，旨在建立“生成、验证、分析、修复、改进、迭代”的闭环改进流程，实现代码质量的持续提升。该机制包含三个关键步骤：错误模式识别、测试反馈分析与针对性修复指导生成。

\subsubsection{问题分析。}
传统代码生成系统采用线性流程，测试失败时仅简单重试，不分析错误根因也不做针对性优化。密码学代码错误涉及算法逻辑、实现细节、环境配置等多维因素，单纯重试无法解决根本问题；现有通用修复策略也缺乏密码学专业支撑，难以适应填充模式调整、初始向量管理等特殊修复需求。

\subsubsection{方法设计。}
自适应改进机制包含三个核心组件：错误模式识别、测试反馈分析与针对性修复。各组件设计要点如下：
\begin{enumerate}
    \item \textbf{错误模式识别} \\
    结合密码学特点，采用多级规则检测与函数化封装，从输出完整性、内容一致性、算法完整性与占位符等维度快速归类，并辅以语法检查，形成综合检测链。

    \item \textbf{测试反馈分析} \\
    依次执行输出缺失、长度差异、明文误输出、导入缺失与错误模式汇总五步，将测试输出压缩为结构化错误特征，供修复指导生成使用。

    \item \textbf{针对性修复指导生成} \\
    依据特征向量生成含错误类型、修复要点与示例代码的结构化指令；采用「模板匹配 + 规则生成」，按算法、语言与错误细节动态填充，兼顾可操作性并降低 LLM 误解。
\end{enumerate}

\subsubsection{闭环改进流程。}
基于上述三个核心组件，本文构建了“生成、验证、分析、修复、改进、迭代”的六步闭环改进机制，具体步骤如图 \ref{fig:loop-process} 所示：
\begin{enumerate}
    \item \textbf{初始生成}。向 LLM 输入分层组合提示，生成初始代码；
    \item \textbf{代码验证}。使用 NIST、FIPS 等标准测试向量进行自动化测试，收集实际输出、错误信息等反馈数据；
    \item \textbf{错误分析}。通过错误模式识别确定错误类型，基于测试反馈分析生成错误特征向量，定位核心错误点；
    \item \textbf{修复指导生成}。根据错误特征向量输出结构化的针对性修复指令；
    \item \textbf{代码改进}。结合修复指令引导 LLM 生成改进代码；
    \item \textbf{迭代终止}。若改进代码通过测试或达到最大迭代次数（3 次）则终止；否则回到第 2 步继续迭代。
\end{enumerate}
\begin{figure}[htbp]
  \centering
  \includegraphics[width=0.85\textwidth]{SixStepIteration2.pdf}
  \caption{六步闭环改进流程}
  \label{fig:loop-process}
\end{figure}
如图\ref{fig:adaptive-iteration}，以 DES-CFB 等模式为例，闭环依次完成初始生成、标准向量验证、错误分析、修复指导注入与代码改进；未通过则重复后三步直至上限 3 轮。
\begin{figure}[htbp]
  \centering
  \includegraphics[width=0.55\textwidth]{adaptive_improvement_iteration.pdf}
  \caption{Complete iteration of the adaptive improvement mechanism for DES-CFB mode code generation}
  \label{fig:adaptive-iteration}
\end{figure}
\subsection{分层可组合的领域提示工程方法}
分层可组合的领域提示工程方法旨在通过设计分层提示模板，对 LLM 实现精准引导。本文提出的提示架构由四层组成：通用模板层、LLM 模板层、语言模板层与算法模板层，每层包含多个精炼子模板。

\subsubsection{分层提示架构设计。}
针对代码跨场景与跨模型兼容性差等问题，我们提出分层提示架构设计。该多层框架将复杂的提示工程需求分解到不同粒度，便于对通用规范、算法特性、模型特点与语言特性进行针对性优化。四层核心结构如下：

\begin{enumerate}
    \item \textbf{通用模板层}。定义系统角色、代码规范、输入/输出格式及基线要求。
    \item \textbf{LLM 模板层}。针对不同 LLM 特性进行优化，按提供商与 \texttt{llm.yaml} 等侧载模板差异化拼装，减少跨场景错误。
    \item \textbf{语言模板层}。按编程语言明确密码库调用规则、代码风格与防错指南，增强代码可编译性。
    \item \textbf{算法模板层}。聚焦算法特有属性（原理、核心步骤、参数要求）引入按算法定制的提示，提升代码验证通过率。


\end{enumerate}

\subsubsection{失败驱动的提示精炼与模式适配。}
在\textbf{四层架构}之上另有二机制：由失败迭代驱动的提示改写，以及按上下文在 \textbf{Full}/\textbf{Compact} 间切换的长度适配。

\textbf{失败驱动的提示精炼。}高 \textbf{FTPR} 依赖跨 \textbf{算法–语言–LLM} 的试错；缺头文件/未定义符号、明文当密文、十六进制宽度、\textbf{IV}/库/\textbf{API} 误用等反复出现。对应约束写入各组合模板（头文件与辅助函数顺序、\texttt{\%02x} 输出、\textbf{CBC} 与 \texttt{TEST\_IV}、库绑定等），验证前可\textbf{轻量修补}。该环节\textbf{最耗时}，对 FTPR 影响最大。

\subsubsection{各阶段的具体挑战与提示对策。}
\label{subsec:challenges_prompt_table}
\textbf{生成}侧固化头文件、自包含、按字节密文、\textbf{CBC} 与输入校验，脆弱组合辅以正/反例；截断时用带 \texttt{last\_error} 的重试。\textbf{改进}侧拼接失败与编译/测试反馈，回注 \texttt{\#include}、\textbf{CBC/IV} 与关键 \textbf{API} 用法。各层将典型修复沉淀为\textbf{可复用约束}。

\subsection{面向本地模型的教师示例注入（检索式蒸馏）}
\label{subsec:distillation}
为提高离线本地模型（如 Qwen2.5-Coder:7B）的 FTPR，我们以云端通过实现为\textbf{教师}，按算法/模式/语言/操作（可并入错误簇）检索注入。改进阶段按模板拼装测试反馈、可选修复摘要、教师块与学生代码，实现\textbf{检索—注入}。

首轮在提示前注入同条件少样本；改进轮可追加同题教师参考（如 SM4-CBC C）。数据来自自动 JSONL 或历史回填，仅本地路径启用。

在 45 个算法–模式–语言组合上，注入前后整体通过率见表~\ref{tab:distillation_overall}（该表为子集统计）。
\begin{table}[H]
  \centering
  \footnotesize
  \caption{教师示例注入前后整体通过率（本地模型；45 项）。}
  \label{tab:distillation_overall}
  \begin{tabular}{l c c}
    \toprule
    & 通过数 & 通过率 \\
    \midrule
    注入前整体通过率 & 37 / 45 & 82.2\% \\
    注入后整体通过率 & 45 / 45 & 100.0\% \\
    提升             & +8 项   & +17.8\% \\
    \bottomrule
  \end{tabular}
\end{table}

\subsection{密码学代码错误模式分类系统}
\subsubsection{错误模式分类。}
密码学代码错误模式分类系统是本文方法的理论基础，旨在为错误识别与纠正提供系统化框架。通过对 3095 个真实测试用例的错误类型分析，本文构建了包含四大类及若干子类的错误分类体系。表 \ref{tab:error_classification} 概括了四类及其代表性示例；更细化的分层约束与失败驱动对策见第~\ref{subsec:challenges_prompt_table}~节及上文。
\begin{table}[htbp]
  \centering
  \footnotesize
  \caption{密码学代码错误模式分类（四类）}
  \label{tab:error_classification}
  \setlength{\tabcolsep}{6pt}
  \begin{tabular}{p{0.3\textwidth}p{0.65\textwidth}}
    \toprule
    大类 & 代表性示例 \\
    \midrule
    算法理解错误 & 模式混淆（如 OFB-64 与 OFB-8）、填充/密钥编排错误。 \\
    实现细节错误 & IV 使用（默认/随机或输出中含 IV）、填充处理、字节序。 \\
    环境配置错误 & 环境变量命名（KEY 与 TEST\_KEY）、缺库、OpenSSL 版本。 \\
    代码结构错误 & \texttt{goto cleanup} 无效、缺头文件、辅助函数不完整（未定义引用）。 \\
    \bottomrule
  \end{tabular}
\end{table}

\subsubsection{修复策略设计。}
针对每类错误，本文设计专用修复策略，保证修复方案的针对性与有效性，核心思路是根据错误类型提供明确的修复要求与示例代码。

算法理解错误修复侧重补充专业知识：密钥长度错误需明确 DES 密钥的标准长度并指导截取前 8 字节；OFB 模式错误需指定使用 OFB64 模式并推荐 EVP\_des\_ofb()；填充模式错误需说明库的自动填充机制并指导删除手动填充代码；密钥编排错误需提供正确置换表的定义与使用示例。

\subsubsection{错误模式匹配算法。}
错误模式匹配采用精确匹配、模式匹配与启发式匹配三种多级策略，协同实现错误模式的快速、精确识别，为生成修复指导提供可靠依据。

\section{实验与结果分析}
本节通过对比实验、消融实验与案例分析验证所提方法的有效性，从多维度系统评估密码学代码生成的性能。

\subsection{实验设置}
\subsubsection{数据集构建。}
\label{subsec:eval_protocol}
为保证实验评估的真实性与代表性，本文构建了包含 3095 个真实测试用例的密码学代码生成数据集，其统计特征见表~\ref{tab:dataset_stats}。这些用例覆盖 4 种主流密码算法与 3 种编程语言（C/C++/Python），均使用 NIST、FIPS 及中国国密标准（GM/T）规定的标准测试向量进行验证。\textbf{除表中另行标注的子集外，后文指标均在同一固定协议下汇总}（一致测试桩、有限次闭环改进等）：便于对齐与复现，\textbf{不}应直接解读为开放部署中的普适表现。


% 数据来源：仓库根目录 code_history.db，表 code_history，条件 test_success=1（与 test_data.yaml 标准向量评测一致的成功入库记录）。
\begin{table}[htbp]
  \centering
  \caption{密码学代码生成数据集统计（3095 条；按算法与语言分布）}
  \label{tab:dataset_stats}
  \setlength{\tabcolsep}{3pt}
  \begin{tabular}{p{0.12\textwidth}
    >{\centering\arraybackslash}p{0.16\textwidth}
    >{\centering\arraybackslash}p{0.30\textwidth}
    >{\centering\arraybackslash}p{0.30\textwidth}
    }
    \toprule
    算法 & 测试用例数 & 语言分布 (C/C++/Python) & 标准测试向量来源 \\
    \midrule
    DES       & 2520 & 697/693/1130 & FIPS 46-3 \\
    AES       & 335  & 102/116/117  & FIPS 197 \\
    SM4       & 184  & 60/63/61     & GM/T 0002-2012 \\
    RSA       & 56   & 19/19/18     & FIPS 186-4 \\
    \midrule
    合计      & 3095 & 878/891/1326 & -- \\
    \bottomrule
  \end{tabular}
\end{table}


\subsubsection{基线系统选择。}
为验证所提方法的有效性与相对优势，本文设置三类对照：\emph{通用外部基线}、\emph{代表性相关工作对比}与\emph{自身模块的消融}。

\textbf{外部基线（通用设置）。}使用相同云端 LLM，仅采用\textbf{单一通用提示}，模拟无任何领域设计的通用代码 LLM 的典型用法。如表~\ref{tab:ablation_prompt} 所示，该基线在\textbf{与下文两张消融表相同}的 DeepSeek、豆包与 Claude ×四算法全支持模式×三语言全网格上仅达 25.0\% FTPR，而四层可组合提示在同一设置下达到 100.0\% FTPR，差距显著，说明分层提示与测试反馈驱动的改进对密码学代码的必要性。

\textbf{相关工作对比。}除通用外部基线外，本文还选取 SecCoder ACL、SVEN LM、Self-Refine 与 AgentCoder 等代表性方法进行对比，并在后文的表~\ref{tab:rw_same_protocol_des} 中给出同一 DES 子集口径下，不同方法在密码学代码生成任务中的实际表现。

\textbf{消融基线}（各自去掉所提方法的一个核心模块以隔离组件贡献）：
\begin{enumerate}[leftmargin=1.2em, labelsep=0.4em]
  \item \textbf{无测试反馈的所提方法}。关闭基于向量测试失败的自改进闭环。 
  \item \textbf{无分层提示的所提方法（\texttt{common\_only}）}。仅保留 \texttt{base\_prompt.yaml} 通用基座，不使用算法/语言等分层模板。
  \item \textbf{无任何提示的所提方法（\texttt{no\_prompt}）}。不加载分层领域模板；用户消息**仅含单行任务描述**（算法、模式、语言与操作），不含通用基座与算法/LLM 层等。
\end{enumerate}

\textbf{消融与错误统计口径。}表~\ref{tab:ablation_main} 与表~\ref{tab:ablation_prompt} 均在\textbf{DeepSeek、Doubao与 Claude}三家云端 LLM 上、于\textbf{四算法全支持工作模式$\times$三语言}的完整任务网格上汇总（VPR/FTPR 为跨三家提供商、该全网格上的聚合结果）。表~\ref{tab:error_repair} 的四类分布与「修复后通过率」则仅基于\textbf{DeepSeek} 在上述\textbf{全算法全模式}网格上写入的失败记录聚合。

\subsubsection{评估指标定义。}
为全面衡量密码学代码生成质量，定义六项评估指标，涵盖正确性、效率与修复能力：
\begin{itemize}[leftmargin=1.2em, labelsep=0.4em]
  \item \textbf{生成成功率（GSR）}：成功生成可运行代码的测试用例比例；
  \item \textbf{验证通过率（VPR）}：通过语法检查与编译的测试用例比例；
  \item \textbf{功能测试通过率（FTPR）}：输出与标准测试向量一致的测试用例比例；
  \item \textbf{错误修复成功率（ERSR）}：改进后通过功能测试的错误用例比例；
  \item \textbf{平均生成时间（AGT）}：每测试用例的平均生成时间（单位：秒）。
\end{itemize}

\subsection{对比实验结果与分析}
\subsubsection{整体性能对比。}
表 \ref{tab:llm_performance} 给出了所提方法在不同 LLM 上的整体性能对比，主要结论如下：

\begin{table}[H]
  \centering
  \footnotesize
  \caption{所提方法在不同 LLM 上的整体性能}
  \label{tab:llm_performance}
  \setlength{\tabcolsep}{4pt}
  \begin{tabular}{
    p{0.10\textwidth}               
    >{\centering\arraybackslash}p{0.19\textwidth}  
    >{\centering\arraybackslash}p{0.20\textwidth}  
    >{\centering\arraybackslash}p{0.06\textwidth}  
    >{\centering\arraybackslash}p{0.06\textwidth}  
    >{\centering\arraybackslash}p{0.07\textwidth} 
    >{\centering\arraybackslash}p{0.07\textwidth}  
    >{\centering\arraybackslash}p{0.08\textwidth} 
  }
    \toprule
    LLM               & 提示长度 (Python/C/C++) & 生成时间 (Python/C/C++) (s) & GSR (\%) & VPR (\%) & FTPR (\%) & ERSR (\%) & AGT (s) \\
    \midrule
    OpenAI            & 17542/28952/22986            & 36.3/60.9/65.8                     & 100.0    & 100.0    & 100.0     & 57.58     & 164.80  \\
    Deepseek          & 8810/17227/13345             & 21.5/53.0/33.2                     & 100.0    & 100.0    & 100.0     & 60.18     & 54.82   \\
    Doubao            & 13801/24722/19702            & 68.0/158.8/94.0                    & 100.0    & 100.0    & 100.0     & 60.56     & 145.33  \\
    Claude            & 8334/15959/12139             & 30.7/65.5/61.6                     & 100.0    & 100.0    & 100.0     & 27.08     & 32.80   \\
    CodeX             & 17542/28952/22986            & 201.8/334.4/275.5                  & 100.0    & 100.0    & 100.0     & 28.57     & 245.09  \\
    qwen2.5-coder:7b  & 14011/22263/17619            & 3182.7/2857.6/2378.4               & 100.0    & 100.0    & 88.89     & 16.67     & 3588.42 \\
    \bottomrule
  \end{tabular}
\end{table}

\begin{figure}[H]
  \centering
  \includegraphics[width=\textwidth]{figures/llm_performance_overview.pdf}
  \caption{各 LLM 性能概览（与表~\ref{tab:llm_performance} 同源数据）。}
  \label{fig:llm_performance_overview}
\end{figure}

\begin{itemize}[leftmargin=1.2em, labelsep=0.4em]
  \item 云端 LLM 在表~\ref{tab:llm_performance} 对应设置下 GSR/VPR 稳定居高，体现分层提示与闭环改进的收益（前提见第~\ref{subsec:eval_protocol}~节）；
  \item 代码提取准确率 100\%，与通用模板层规定的输出格式一致；
  \item 首次生成即通过验证的比例超过 90\%；
  \item \textbf{qwen2.5-coder:7b（离线）}同网格下 GSR/VPR 为 100\%，FTPR 88.89\%；AGT 约 3588\,s，显著高于云端，适合隐私优先、可容忍高时延的场景。
\end{itemize}

\subsubsection{相关工作对比。}
为更直接比较本文方法与代表性相关工作的表现，我们在\textbf{同一 DES 子集评测设置}（ECB/CBC/CFB/OFB × Python/C/C++，共 12 格）下整理并报告与各基线对齐的 GSR/VPR/FTPR。表~\ref{tab:rw_same_protocol_des} 汇总主要对比：各基线在同设置下的 FTPR 显著低于本文 DeepSeek 结果。

\begin{table}[H]
  \centering
  \footnotesize
  \caption{同设置相关工作对比（DES 全模式 × DeepSeek）}
  \label{tab:rw_same_protocol_des}
  \setlength{\tabcolsep}{4pt}
  \begin{tabular}{p{0.42\textwidth} c c c}
    \toprule
    基线 & GSR (\%) & VPR (\%) & FTPR (\%) \\
    \midrule
    SVEN & 100.0 & 0.0 & 0.0 \\
    Self-Refine+DeepSeek & 100.0 & 83.33 & 58.33 \\
    SecCoder+DeepSeek & 100.0 & 66.67 & 66.67 \\
    AgentCoder & 100.0 & 41.67 & 0.0 \\
    本文方法+DeepSeek & 100.0 & 100.0 & 100.0 \\
    \bottomrule
  \end{tabular}
\end{table}



\subsubsection{代码生成与提示优化中的挑战。}
\label{subsec:challenges_prompt}
\textbf{维持高 FTPR} 依赖反复失败上的提示迭代，是实验中最耗时、难度最高的环节；写法与对策已概括于第~\ref{subsec:challenges_prompt_table}~节。

常见失败包括缺 \texttt{stdint.h}/未定义符号、明文当密文、\texttt{\%04x} 与按字节 \texttt{\%02x} 不一致、缓冲与输入校验问题等。\textbf{生成}侧强化约定与带 \texttt{last\_error} 的重试；\textbf{改进}侧优先消除编译错误并回注 API 与控制流要点，必要时预插头文件。


\subsubsection{错误模式分类效果。}
对 DeepSeek 在全算法全模式上写入的功能测试失败记录经启发式归入四类；共 101 条失败样本，合计修复后通过率为 87.13\%（见表~\ref{tab:error_repair}）。实现细节类记录最多；算法理解类修复后通过率相对较低；环境配置与代码结构类在日志记录层面上均已跑通。主要结论如下：
\begin{enumerate}[leftmargin=1.2em, labelsep=0.4em]
  \item \textbf{分布与难点。}实现细节类失败条数最多（41，40.59\%）；算法理解类修复后再通过最低（77.14\%）。
  \item \textbf{易收敛类型。}环境配置类与代码结构类在观测口径下修复后通过率均为 100.0\%，特征相对明确（如环境变量名、缺头文件、结构问题），针对性修复指导易见效。
  \item \textbf{整体效果。}自 0\% 基线至合计 87.13\% 的修复后通过率，验证了错误模式分类与针对性修复的有效性。
\end{enumerate}

\begin{table}[H]
  \centering
  \footnotesize
  \caption{不同错误类型的分布与修复效果（DeepSeek）}
  \label{tab:error_repair}
  \setlength{\tabcolsep}{3pt}
  \begin{tabular}{
    p{0.28\textwidth}                       
    >{\centering\arraybackslash}p{0.08\textwidth} 
    >{\centering\arraybackslash}p{0.12\textwidth} 
    >{\centering\arraybackslash}p{0.16\textwidth}  
    >{\centering\arraybackslash}p{0.18\textwidth}  
  }
    \toprule
    错误类型               & 数量 & 占比 & 修复前通过率 & 修复后通过率 \\
    \midrule
    算法理解  & 35    & 34.65\%    & 0.0\%                & 77.14\%               \\
    实现细节  & 41    & 40.59\%    & 0.0\%                & 87.80\%               \\
    环境配置  & 12    & 11.88\%    & 0.0\%                & 100.0\%               \\
    代码结构  & 13    & 12.87\%    & 0.0\%                & 100.0\%               \\
    \midrule
    合计      & 101   & 100.0\%    & 0.0\%                & 87.13\%               \\
    \bottomrule
  \end{tabular}
\end{table}

\subsection{消融实验}
通过消融实验量化所提方法各核心模块的贡献。表~\ref{tab:ablation_main} 给出主要消融实验结果，关键结论如下。

\begin{table}[H]
  \centering
  \footnotesize
  \caption{主要消融实验结果（DeepSeek、豆包与 Claude；四算法全支持模式$\times$三语言）}
  \label{tab:ablation_main}
  \setlength{\tabcolsep}{4pt}
  \begin{tabular}{l c c c}
    \toprule
    系统配置                  & VPR (\%) & FTPR (\%) & 性能下降 (VPR/FTPR) \\
    \midrule
    完整所提方法                  & 100.0    & 100.0     & -- / --                     \\
    无测试反馈改进         & 100.0    & 77.08     & 0.0\% / -22.92\%            \\
    无分层提示架构（仅 base_prompt.yaml）  & 60.42    & 50.0      & -39.58\% / -50.0\%          \\
    无任何提示  & 0.0    & 0.0     & -100.0\% / -100.0\%         \\
    \bottomrule
  \end{tabular}
\end{table}

\begin{enumerate}[leftmargin=1.2em, labelsep=0.4em]
  \item \textbf{测试反馈改进主要作用于 FTPR}：去除闭环后 FTPR 降至 77.08\%，说明\textbf{向量一致性功能正确性仍依赖测试反馈}。
  \item \textbf{分层提示架构是功能正确性的关键}：仅 \texttt{common\_only} 时 VPR/FTPR 为 60.42\%/50.0\%。
  \item \textbf{完全撤除领域栈则指标塌缩}：\texttt{no\_prompt} 下仅保留单行任务描述，表中 VPR/FTPR 仍为 0\%。
\end{enumerate}


\subsubsection{提示优化消融。}
针对提示策略的进一步消融验证了分层可组合设计的价值。表~\ref{tab:ablation_prompt} 给出不同提示策略的对比结果：
\begin{itemize}[leftmargin=1.2em, labelsep=0.4em]
  \item 随\textbf{提示模板层次}由通用向「通用+LLM+语言」递进，VPR/FTPR 逐步升高（表~\ref{tab:ablation_prompt}），说明组合多层模板能充分发挥领域知识与模型特性；
  \item \textbf{四层可组合提示策略（所提方法）}达到最佳性能（FTPR 100.0\%），显著优于单一通用模板（25.0\%）及中间档位（79.17\%--91.67\%），证实分层提示工程能有效弥补通用 LLM 在密码学知识上的不足。
\end{itemize}

\begin{table}[H]
  \centering
  %\footnotesize
  \caption{提示策略消融实验结果（DeepSeek、豆包与 Claude；四算法全支持模式$\times$三语言）}
  \label{tab:ablation_prompt}
  \setlength{\tabcolsep}{4pt}
  \begin{tabular}{l c c c}
    \toprule
    提示策略               & GSR (\%) & VPR (\%) & FTPR (\%) \\
    \midrule
    单一通用模板         & 100.0    & 45.83    & 25.0      \\
    通用+LLM              & 100.0    & 87.5     & 79.17     \\
    通用+LLM+语言        & 100.0    & 91.67    & 91.67     \\
    通用 +LLM+ 语言+算法（本文方法） & 100.0 & 100.0 & 100.0 \\
    \bottomrule
  \end{tabular}
\end{table}

\section{结论与未来工作}
\subsection{研究结论}
针对通用 LLM 在密码学代码上错误率高、领域适配不足的问题，本文将测试反馈闭环、四层可组合提示、\textbf{失败驱动的提示精炼}与错误模式分类整合为可复现流程。

在 3095 例网格与第~\ref{subsec:eval_protocol}~节所述协议下，生成/提取成功率达 100\%；完整方法 FTPR 由 45.1\% 升至 93.33\%，优于单一通用提示基线。

贡献可概括为三点：\emph{框架}上统一错误分类、闭环与分层提示；\emph{技术}上高成功率提取与失败驱动精炼相配合（表~\ref{tab:error_repair} 修复后再通过约 87.13\%）；\emph{工程}上 101 份任务模板与多后端脚本便于复现与小规模试用。

\subsection{研究局限}
错误模式库由 3095 条归纳；表~\ref{tab:error_repair} 仅来自 DeepSeek 全网格失败日志。算法理解类修复后再通过约 77.14\%，复杂逻辑仍是短板。

覆盖限于 DES/AES/SM4/RSA；评测以向量一致性为主，未系统纳入侧信道等，不宜外推为线上安全保证。

\subsection{未来工作方向}
拟扩充真实失败样本与提示资产；探索由用户反馈轻量更新模板与评测配置。

% \section*{Acknowledgements}
% This work was supported by Sichuan University. The authors thank the teachers and classmates of the School of Cyber Science and Engineering of Sichuan University for their valuable suggestions and discussions. Doubao was used to produce parts of the work, including text, tables, charts, codes, data, and quotations.

\begin{credits}
\subsubsection{\ackname}
% 终稿可在此写明资助与致谢（与投稿须知一致）。
\subsubsection{\discintname}
作者声明不存在与本研究内容直接相关的竞争性利益。
\end{credits}

\renewcommand{\refname}{参考文献}
\renewcommand{\bibname}{参考文献}
\bibliographystyle{splncs04}
\bibliography{refs} 
%% Appendix (if needed)
%\appendix
%\section{Appendix Title}
%Appendix content here.

\end{document}